\section{Limitations and Conclusion}
\label{sec:conclusion}

\noindent\textbf{Limitations.}
\ours assumes depth and odometry. Its primary EB-Navigation condition receives simulator depth and exact pose,
whereas the benchmark-prompt baseline receives RGB only; therefore their gap cannot be assigned solely to the
planning algorithm. Estimated depth reduces SR to 45.7\%, and we do not remove the sensing assumptions.
EB-Navigation ends an episode automatically on success; the exhaust-on-arrival rule uses that signal, so we report
its ablation and evaluate on RoboTHOR, where the agent must decide to stop.
Appearance-only goals remain hard because a 4B model's candidate set often misses the target. Single-room scenes
and same-category look-alikes limit how far the EB-Navigation results generalize.

\noindent\textbf{Conclusion.}
Small VLMs can provide semantic cues that a geometric controller can use for navigation. In the present artifacts,
semantic-to-metric grounding reaches a three-run mean of 79.6\% on EB-Navigation and 59.6\% SR on RoboTHOR;
independent, input-matched baselines remain necessary before making a causal performance claim.
